\section{My Evaluation}
Footnote\footnote{Author, \emph{Title} (City: Publisher, Year), page.}

0. Intruduction

First of all, the model sticked to most of the rules I have given to it;
it sticked to $800$ word limit ($787$ to be exact) and followed the given structure in the prompt.

But it broke the quotation rule, as {\ttfamily\small "undreamed of before,"} is the only direct quote from the source, although the prompt asked for at least five.
And even there, it did not get the punctuation right.
Being unable to quote correctly, even from just the prompt,
is probably the result of the model being only $9$ billion parameters and quantised to fit in under \qty{8}{\giga\byte}.

1. What kind of source is it?
- good: author

The interpretation of the article is mostly accurate,
for example
But it made some mistakes in the details.

A significant difficulty that arose from the prompt was that I told it to play the role of \enquote{an experienced historian of modern Europe},
while demanding to flag any information (a date, a name, an event, an outcome) it was unsure about.
That can be interpreted as a contradiction, because an experienced historian can usually look these things up, but the model could not do that.
Therefore, it confidently made up some information where it should have flagged it instead, probably to seem more like a real historian.
An example would be
- quote sources and prove model wrong

4. Why does the source matter?
- add a significant point the model missed

All in all, small open models are already very useful for basic tasks such as summarisation.
But especially for a reliable and deep primary source analysis, they are not yet at the level of bigger models, and certainly not human historians.
Human historians can look up information and verify it, therefore handeling a way bigger amount of input, which is fundamental for a primary source analysis.
This might get better in the future, but the fundamental restriction comes from the underlying model architecture.
$9$ billion parameters is simply not enough to reliably handle such a large input.
But there are already efforts to address this problem, for example the Per-Layer Embedding (PLE) technology,
introduced by Google's \href{https://ollama.com/library/gemma4}{gemma4} model.
